% ================= 1. INTRODUCTION =================
\section{Introduction}

\subsection{Background \& Problem Statement}
\label{sec:background}

The Cubli is a reaction-wheel inverted pendulum: a cube-shaped platform that balances on an edge or a corner using only internal actuation, with no external supports or contact forces beyond the single line or point of contact with the ground. The concept was introduced by the Institute for Dynamic Systems and Control (IDSC) at ETH Zurich \cite{gajamohan2012cubli, gajamohan2013cubli}, whose original prototype ($15\times15\times15$~cm) mounts one brushless DC motor and flywheel on each of three orthogonal faces. Commanding a change in wheel speed produces, by conservation of angular momentum, an equal and opposite reaction torque on the cube body; braking a wheel abruptly transfers its stored angular momentum to the body in a near-impulsive event, which is what allows the cube to \emph{jump up} from resting on a face into a balance on an edge or corner.

Dynamically, the system is an \emph{underactuated}, nonlinear, unstable-equilibrium control problem: three actuators (the wheels) must stabilize a body with more degrees of freedom than actuated inputs, about equilibria (edge and corner balance) that are open-loop unstable and increasingly coupled across axes as the contact geometry shrinks from a line (edge) to a point (corner). This is structurally the same problem as the classical inverted pendulum, extended from a single rotational axis to a fully three-dimensional, momentum-exchange-actuated body — which is why the Cubli has become a standard benchmark platform in the controls community over the past decade.

The problem this project addresses is therefore twofold: (i) derive and validate a dynamic model and controller capable of stabilizing this underactuated system about its unstable equilibria and driving it through commanded reorientations, and (ii) realize this in hardware as a modular, largely parametric design, since final actuator and sensor specifications are not yet confirmed at the time of writing. The target capability set, in order of increasing difficulty, is edge balance, corner balance, controlled rotation, and walking (Section~\ref{sec:objectives}).

\paragraph{Why this is a hard build, not a decorative demo.} It is worth stating plainly why a machine that looks like a cube with three flywheels resists casual replication. The difficulty is not that any single subsystem is exotic --- none is --- but that every subsystem sits close to a physical bound, and the bounds are \emph{coupled}: relieving one tightens another. The following failure modes recur across published and amateur replications, and each is a consequence of that coupling rather than of poor workmanship.

\begin{itemize}
    \item \textbf{The jump-up is an energy problem with no margin.} The wheel must store, at its maximum speed, enough angular momentum to lift the cube's center of mass over the corner. The required per-wheel inertia scales as $M L^{3/2}$ (Section~\ref{sec:jumpup_target}), so mass added \emph{anywhere} --- battery, frame, fasteners --- demands a larger wheel, and the wheel is itself part of $M$. The sizing problem is therefore a fixed-point iteration rather than a one-pass calculation, and it does not always converge: for the present cube the three-wheel set alone accounts for roughly a quarter of the total mass budget, and shrinking the wheel diameter to save space makes the situation worse, not better, because inertia must then be recovered with ballast at a smaller radius.
    \item \textbf{Braking torque fights the structure.} The momentum accumulated over seconds of spin-up must be shed in tens of milliseconds. That impulse passes through one motor shaft, one hub, and one motor mount --- the very parts the mass budget wants thin. Many builds fail here first, and they fail in a characteristic way: the electronics are fine, the controller is sound, and the coupling, hub or mount yields on the first \emph{successful} jump. The structural load case is set by the maneuver that proves the machine works.
    \item \textbf{Sensing, not the control law, sets the balancing limit.} The attitude estimate degrades from two directions at once: MEMS gyroscope bias drifts, and the accelerometer --- the only absolute reference for tilt --- is corrupted by vibration that the controller itself generates by spinning the wheels. Structural design and state estimation are consequently not separable concerns. A wheel half a gram out of balance at rim radius can render a perfectly correct estimator unusable, which is why wheel balance is treated here as an estimation requirement and not merely as a mechanical nicety.
    \item \textbf{The plant is genuinely nonlinear and constrained.} Corner balancing does not decouple into three independent single-axis problems: the axes are kinematically coupled through the contact point and the body inertia tensor. The actuators saturate in \emph{both} torque and speed, and under any persistent disturbance the wheels drift toward their speed limit, so stored momentum must be actively managed rather than left to accumulate. Linearizing about the upright equilibrium is entirely adequate for a balancing demonstration and inadequate for anything more interesting \cite{muehlebach2017cubli}.
\end{itemize}

The common thread is worth making explicit, because it determines how this project is sequenced. Replications typically fail not because the control law was wrong, but because the machine was never capable of the maneuver being asked of it: insufficient wheel inertia, motors that cannot reach their rated speed under load, a brake too slow to approximate an impulse, or a mass budget that quietly grew past the point where the stored energy suffices. Every one of these is decidable on paper, before any material is cut --- and commonly is not decided there. That is the reason sizing is treated in this document as the first engineering activity rather than a detail to be settled once parts arrive.

\subsection{Purpose, Application \& Mission Context}
\label{sec:application}

The Cubli is not presented here as a balancing novelty but as a terrestrial, gravity-loaded demonstrator of \textbf{momentum-exchange attitude control} — the same physical mechanism that governs how real spacecraft point themselves. Nearly every three-axis-stabilized satellite holds and changes its attitude using reaction wheels: flywheels whose commanded speed change produces a reaction torque on the spacecraft body, exactly as in the Cubli. Spacecraft typically carry three or four such wheels (often in a pyramidal arrangement for redundancy) — directly analogous to the Cubli's three orthogonal wheels — and this reaction-wheel approach is the standard for CubeSats and small satellites in particular, making the Cubli representative of a real, currently-flown hardware class rather than an idealized toy problem. Reaction wheels are distinct from control moment gyros (CMGs), which gimbal a constantly-spinning wheel instead of changing its speed to produce much larger torques \cite{votel2012cmgvsrw}; positioning the Cubli explicitly as a \emph{reaction-wheel}, not CMG, demonstrator keeps the application claim precise.

Framed as an ADCS problem, the Cubli's core behaviors map directly onto a spacecraft's day-to-day pointing task:

\begin{itemize}
    \item \textbf{Edge / corner balance} $\rightarrow$ regulation to a commanded reference attitude and disturbance rejection — holding a pointing setpoint against perturbations.
    \item \textbf{Controlled rotation} $\rightarrow$ a slew maneuver: reference tracking to a new commanded attitude, exactly as a satellite re-aims a payload or antenna.
    \item \textbf{Jump / walk} $\rightarrow$ momentum-exchange hopping mobility, flight-proven on small bodies where wheeled traction is impossible. JAXA's Hayabusa2 mission deployed the MINERVA-II1 rovers \cite{yoshimitsu1999minerva} and the MASCOT lander \cite{ho2017mascot} on asteroid Ryugu (2018), both of which move by spinning up and braking an internal flywheel or eccentric mass to hop across the surface, with MASCOT additionally self-righting the same way. Walking, as a sequence of controlled, directional hops, is a genuine step beyond a single ballistic hop and is an active research frontier (e.g.\ ETH's SpaceHopper \cite{spiridonov2024spacehopper}).
\end{itemize}

This gives the project two complementary application narratives rather than one: it is simultaneously a low-cost, hardware-in-the-loop testbed for standard spacecraft \emph{attitude determination and control} (ADCS), and a ground analog for the momentum-exchange \emph{surface mobility} now used on small-body missions. Framed against a companion star-tracker project, which answers \emph{``where am I pointing?''} (attitude determination), the Cubli answers \emph{``drive me to that reference and hold it''} (attitude control) — the two halves of a complete ADCS loop.

\paragraph{Why the problem is still open.} The nominal Cubli problem --- balance on a corner, jump up --- is solved, and was solved a decade ago on well-instrumented, precisely manufactured, generously funded hardware \cite{gajamohan2012cubli, gajamohan2013cubli, muehlebach2017cubli}. Reproducing that result is an engineering exercise, not a research contribution, and this document does not claim otherwise. The interesting questions begin immediately \emph{after} that point, and they are not Cubli-specific: they are the questions a spacecraft attitude-control engineer faces, in a form small enough to be answered on a bench.

\begin{itemize}
    \item \textbf{Control under simultaneous torque and momentum saturation.} A jump saturates the wheels by construction --- the maneuver is defined by running them to their speed limit and dumping the momentum. Steering a nonlinear attitude system to a target while respecting a finite momentum envelope, and desaturating without losing the pointing objective, is still handled largely case-by-case with tuned heuristics rather than by a general method.
    \item \textbf{Recovery from large attitude excursions.} Linear controllers valid near upright say nothing useful once the body is far from equilibrium, which is precisely the regime a tumbling or mis-deployed spacecraft occupies. Nonlinear and predictive methods make claims there, and verifying those claims empirically requires a system in which a failed recovery costs a re-print rather than a mission.
    \item \textbf{Robustness to badly known parameters.} Real inertia tensors, friction coefficients and center-of-mass offsets are known to a few percent at best, and shift with assembly, thermal state and wear. Adaptive and robust designs are overwhelmingly validated against \emph{simulated} parameter errors, which are drawn from the same distribution the designer assumed.
    \item \textbf{Whether high-performance methods survive commodity sensing.} Attitude-control techniques developed against clean measurements must eventually run on a MEMS IMU with drifting bias and vibration-corrupted accelerometers. Whether a given method degrades gracefully or collapses under that sensing is an empirical question, and it is settled by hardware rather than by analysis.
\end{itemize}

\paragraph{Why a Cubli, and why built from scrap.} The Cubli is the cheapest honest test of the four questions above. Simulation cannot answer them, because the effects that decide the outcome are exactly the ones a model omits: unmodeled friction, frame resonance excited by the wheels, and the interaction between wheel imbalance and estimator bias. These are what separate a controller that works in simulation from one that works on hardware. At the other extreme, a full-scale testbed --- air bearing, vacuum chamber, flight-grade wheels --- costs orders of magnitude more and, decisively, \emph{cannot be crashed repeatedly}. The scrap-built cube occupies the useful middle: physical enough to be honest, cheap enough to be abused. The choice to build it from commodity and salvaged parts is therefore methodological rather than merely budgetary, on four counts:

\begin{itemize}
    \item \textbf{Cheap failure means the aggressive experiments actually get run.} A controller that must never fail cannot be tested at its limits, and a testbed too valuable to break silently biases the experimental program toward maneuvers already known to work.
    \item \textbf{Cheap parts are a stronger test than good parts.} Loose tolerances, an imperfectly balanced wheel and a drifting IMU are not defects to be apologized for --- they \emph{are} the parameter uncertainty and disturbance content that robust and adaptive methods claim to tolerate. A demonstration on deliberately imperfect hardware therefore carries evidence that a demonstration on precision hardware does not.
    \item \textbf{Constraints force honest sizing.} With a fixed motor, a fixed print volume and a fixed budget, a marginal design cannot be rescued by specifying a larger component. This is the entire justification for the hex-nut-in-pocket ballast scheme of Section~\ref{sec:rw_sizing}: it buys inertia at the outer radius using standard steel nuts and no machined part, and whether it closes is a calculation that must be completed \emph{before} printing, not discovered afterwards.
    \item \textbf{A commodity bill of materials is reproducible.} Every part is orderable or salvageable by any university group, so the result can be checked by others --- which is the difference between a demonstration and a claim.
\end{itemize}

\paragraph{Mission context.} The subsystem exercised here is the one a spacecraft uses for pointing: momentum-exchange attitude control under actuator saturation, with desaturation, on an underactuated body whose inertia is imperfectly known. In one respect the terrestrial case is strictly \emph{harder} than the orbital one. In orbit, a marginal attitude controller degrades slowly and often invisibly --- pointing error grows, wheels creep toward saturation, and the fault may take orbits to become obvious. Under gravity, the destabilizing torque never vanishes and never relents: an inadequate controller falls over immediately and unambiguously. The demonstration is, in that sense, self-evaluating, and this is a feature of the platform rather than a limitation of it.

\subsection{Objectives}
\label{sec:objectives}

The project targets four behaviors, deliberately ordered as a difficulty ramp in which each milestone is a control precondition for the next:

\begin{enumerate}
    \item \textbf{Edge balance} — stabilize the cube on a single edge, the least unstable equilibrium (contact is a line, one primary axis of instability).
    \item \textbf{Corner balance} — stabilize the cube on a single vertex, the fully unstable equilibrium (point contact, all three axes coupled and unstable simultaneously).
    \item \textbf{Controlled rotation} — execute a commanded slew between equilibria (e.g.\ edge to edge, or in-place reorientation), demonstrating reference tracking rather than pure regulation.
    \item \textbf{Walking} — chain repeated, directional controlled falls between edges/corners into a net translation across a surface, re-stabilizing after each transition.
\end{enumerate}

These map one-to-one onto the applications of Section~\ref{sec:application}: (1)–(2) demonstrate pointing regulation and disturbance rejection of increasing difficulty; (3) demonstrates slew/reference-tracking; (4) demonstrates momentum-exchange mobility of the kind flown on Hayabusa2. Milestones (1)–(3) are treated as committed deliverables for this design; walking (4) is scoped as the stretch objective, contingent on hardware confirmation and the jump-up feasibility check (angular momentum required to tip the cube vs.\ available motor/wheel torque), to be resolved once final components are known.